\section{Marco teórico}
\subsection{Pronósticos}
\label{sec:marco-pronosticos}

Una serie de tiempo es una sucesión de observaciones $y_t$ de una variable registradas en períodos consecutivos e igualmente espaciados.
Su comportamiento suele describirse mediante un nivel, una tendencia que lo desplaza con el tiempo, una estacionalidad que repite un patrón con período fijo $m$ y un componente aleatorio sin estructura.
Pronosticar consiste en estimar valores futuros $\hat{y}_{T+h}$ a partir de las observaciones disponibles hasta el período $T$, donde $h$ es el horizonte del pronóstico \cite{hyndman_athanasopoulos_fpp3}.

\subsubsection{Promedio móvil}
La técnica del promedio móvil supone que la serie fluctúa alrededor de un nivel constante, $y_t = b + \varepsilon_t$, con un ruido $\varepsilon_t$ de media cero y varianza constante \cite{taha_io_2004}.
El nivel se estima con el promedio de las $n$ observaciones más recientes, a las que se asigna igual importancia:
\begin{equation}
  \hat{y}_{t+1} = \frac{1}{n} \sum_{i=0}^{n-1} y_{t-i}.
\end{equation}
No existe una regla exacta para elegir $n$.
Un valor grande suaviza mejor el ruido cuando la serie es estable, pero suprime la tendencia si los datos crecen o decrecen, mientras que un valor pequeño sigue más rápido los cambios a costa de un pronóstico más ruidoso \cite{taha_io_2004}.

\subsubsection{Métodos estacionales de referencia}
Cuando la serie tiene estacionalidad, el promedio simple de los últimos períodos mezcla momentos del ciclo que no son comparables.
El método ingenuo estacional evita ese problema pronosticando cada período con el valor observado en el mismo momento del ciclo anterior \cite{hyndman_athanasopoulos_fpp3}:
\begin{equation}
  \hat{y}_{T+h} = y_{T+h-m}, \qquad 1 \le h \le m.
\end{equation}
Es el punto de comparación mínimo que un método más elaborado debería superar.

En este trabajo la serie es la cantidad de pedidos por hora de atención, y su estacionalidad combina un patrón dentro del día con otro entre días de la semana.
Ambos se capturan comparando cada hora con la misma hora del mismo día de semanas anteriores.
Si $y_{d,h}$ es la cantidad de pedidos del día $d$ en la hora $h$, el ingenuo estacional es $\hat{y}_{d,h} = y_{d-7,h}$.
La media estacional aplica el promedio móvil a la subserie de cada hora y día de la semana, promediando las últimas $k$ semanas:
\begin{equation}
  \hat{y}_{d,h} = \frac{1}{k} \sum_{j=1}^{k} y_{d-7j,\,h}.
  \label{eq:media-estacional}
\end{equation}
El ingenuo estacional es el caso particular $k = 1$.
Como el rezago mínimo es de una semana, ambos métodos permiten pronosticar una semana completa usando solo datos anteriores a ella.

\subsubsection{Evaluación del error}
Un pronóstico se evalúa con datos que no se usaron para construirlo.
En series de tiempo la partición respeta el orden cronológico: los primeros períodos sirven para ajustar y elegir el método, y un tramo final reservado mide el error que cabría esperar al pronosticar \cite{hyndman_athanasopoulos_fpp3}.
Si además hay que elegir un parámetro, como $k$, se reserva un tramo de validación previo al de prueba, de modo que el error informado no esté influido por la elección.

Con el error $e_t = \hat{y}_t - y_t$ definido como pronóstico menos valor real, las medidas usadas son:
\begin{align}
  \mathrm{MAE} &= \frac{1}{N} \sum_{t} \lvert e_t \rvert, &
  \mathrm{RMSE} &= \sqrt{\frac{1}{N} \sum_{t} e_t^2}, &
  \mathrm{Sesgo} &= \frac{1}{N} \sum_{t} e_t .
\end{align}
El error absoluto medio (MAE) se expresa en las mismas unidades que la serie y es fácil de interpretar.
La raíz del error cuadrático medio (RMSE) penaliza más los errores grandes.
El sesgo indica si el método sobrestima ($\mathrm{Sesgo} > 0$) o subestima ($\mathrm{Sesgo} < 0$) en promedio, algo relevante para dimensionar personal: subestimar la demanda conduce a dotaciones insuficientes.

Cuando las llegadas siguen un proceso de Poisson, incluso un pronóstico que conociera la tasa verdadera $\lambda$ de cada hora tendría error, porque la cantidad observada varía con varianza $\lambda$.
Ese error irreducible es un piso con el que conviene comparar los métodos para saber cuánto margen de mejora queda.
\subsection{Teoría de colas}
Una fila FIFO asigna el siguiente preparador libre al pedido que lleva más tiempo esperando.
El modelo M/M/$c$ supone llegadas Poisson con tasa constante $\lambda$, servicios exponenciales independientes de tasa $\mu$ por preparador, $c$ preparadores equivalentes y espacio de espera ilimitado \cite{bertsekas_colas,taha_io_2004}.
Si las tasas se expresan por hora, la carga ofrecida $a=\lambda/\mu$ mide cuántos preparadores ocuparía la demanda en promedio y $\rho=a/c$ es la fracción ocupada de la capacidad.
Existe equilibrio estacionario cuando $\rho<1$.
La fórmula de Erlang C da la probabilidad de que un pedido deba esperar:
\[
 P_0=\left(\sum_{k=0}^{c-1}\frac{a^k}{k!}+
              \frac{a^c}{c!(1-\rho)}\right)^{-1},
 \qquad P_{\mathrm{espera}}=P_0\,\frac{a^c}{c!(1-\rho)}.
\]
En equilibrio, la espera media es $W_q=P_{\mathrm{espera}}/(c\mu-\lambda)$ y el tiempo medio hasta estar listo es $W=W_q+1/\mu$, ambos en horas.
Por la ley de Little, la cantidad media en espera es $L_q=\lambda W_q$ y la cantidad media en el sistema es $L=\lambda W$.
La comparación con tiempos en minutos requiere convertir la tasa a pedidos por minuto.
\subsection{Simulación de eventos discretos}
La simulación avanza el reloj entre llegadas, inicios y fines de preparación.
En cada evento se actualiza quién espera y quién ocupa un preparador.
Este mecanismo permite que cambie la demanda al pasar de una hora a otra mientras los pedidos ya admitidos conservan su lugar.
En SimPy, cada pedido se representa mediante un proceso que solicita y luego libera una unidad de un recurso compartido \cite{simpy_shared_resources}.
